\noindent\textbf{ZUR NETZAUSGLEICHUNG.}\hfill319
\begin{center}\textbf{[10.]}\par\medskip
\textbf{[Briefwechsel zur Netzausgleichung.]}\end{center}

\begin{center}\textsc{Gauss} an \textsc{Olbers}.\quad Göttingen, 2. November 1823.\end{center}

.\ .\ .\ .\ .\ .\ Ich habe nun die mühsame Ausgleichung aller meiner
Messungen von 1821---1823 abgeschlossen, soweit sie die Hauptdreiecke betrifft.
Nun stimmen sowohl die Summen der Winkel der einzelnen Dreiecke als auch die
Verhältnisse der Seiten in den gekreuzten Vierecken und Fünfecken genau überein,
und zwar nach den strengen Prinzipien der Wahrscheinlichkeitsrechnung, sine
ira et studio, und ohne jede Willkür. Letzteres scheint sich von den Messungen
anderer Geodäten nicht sagen zu lassen. Im ganzen System gibt es 76 Richtungen,
d.\ h.\ 38, jeweils hin und zurück. Bei \textls[120]{keiner} von ihnen hat
die Ausgleichung $1^{\prime\prime}$ betragen. Die größte ist $0^{\prime\prime}{,}813$
bei der Richtung von Nindorf nach Hamburg, wo es bei der rauchartigen Atmosphäre
und der Phasenstörung immer sehr schwierig war, den Michaelisthurm anzuvisieren.
Die nächstgrößte ist $0^{\prime\prime}{,}788$ bei der Richtung von Lüneburg nach
Wilsede. Dort war der Zielpunkt zwar ein Heliotroplicht, aber die Aufstellung
auf einem hölzernen Stativ in der Turmlaterne war gewiss weniger solide als
auf Steinpostamenten, und das Gewicht des Beobachters mag je nach seiner Stellung
Einfluss auf das Instrument gehabt haben. Der \textls[120]{mittlere} Fehler
aller Richtungen im Sinne meiner Theoria Combinationis ist $0^{\prime\prime}{,}486$.

Im ganzen System ergeben sich 26 Dreiecke, in denen ich alle Winkel
gemessen habe. Der größte Fehler in der Winkelsumme beträgt jetzt $2^{\prime\prime}{,}175$
im Dreieck Nindorf-Hamburg-Timpenberg, wofür vor allem die Richtung von
Nindorf nach Hamburg verantwortlich sein dürfte. Der nächstgrößte Fehler
liegt, wie schon oben erwähnt, im Dreieck Brocken-Hohehagen-Hils. Er beträgt
$1^{\prime\prime}{,}806$, wofür jetzt noch vor allem die Richtung vom Hils zum
Brocken verantwortlich sein dürfte.
\par\medskip
.\ .\ .\ .\ .\ .

\clearpage
\noindent320\hfill\textbf{BRIEFWECHSEL.}
\begin{center}\textsc{Gauss} an \textsc{Olbers}.\quad Zeven, 6. Juli 1824.\end{center}

\begingroup
\makeatletter
\renewcommand{\@makefnmark}{\hbox{\normalfont\textsuperscript{*})}}
\makeatother
.\ .\ .\ .\ .\ .\ Die neuen Messungen aus dem Jahr 1823 [haben] im ganzen
System einige, wenn auch sehr kleine Änderungen bewirkt\footnote{Auch habe ich meiner jetzigen Rechnung \textsc{Walbecks} Abplattung \textls[120]{genau} zugrunde gelegt, während es bei der vorigen, wie Sie wissen, durch ein Versehen zu einer, wenn auch äußerst kleinen, Abweichung kam.}.
\par\medskip
.\ .\ .\ .\ .\ .
\endgroup
